\section{Exact bounds for saturated blocks}
\label{sec:second-order-saturated-blocks}

All multigraphs in this section are finite and loopless.
Parallel edges are allowed and are distinguished when paths are counted.
The four-terminal classification in Proposition~\ref{prop:saturated-four-pole}
also holds for such multigraphs.
Its proof forces a spanning pair of paths and a matching between them.
The resulting permutation ladder is simple.

For positive integers $h$, define
\[
 M(2j)=3\cdot2^j-2,\qquad
 M(2j+1)=4\cdot2^j-2.
\]
Thus $M(1)=2$, and $M(h)\le3\cdot2^{h/2}$.

\begin{theorem}\label{thm:saturated-three-pole}
Let $T$ be connected, with three distinct degree-two terminals $a,b,c$.
Every other vertex has degree three.
Suppose that $\beta(T)=h$ and
$P_{ac}(T)=P_{bc}(T)=2^h$.
Then
\[
 P_{ab}(T)\le M(h).
\]
For every $h\ge1$, a simple graph attains equality.
\end{theorem}

\begin{proof}
Saturation implies that every $a$--$c$ path and every $b$--$c$ path meets every cycle.
Indeed, the corresponding affine parity class has $2^h$ members.
Every member is one terminal path together with vertex-disjoint cycles.

We first reduce bridges by induction on $h$.
Every component obtained by deleting a bridge contains a cycle.
To see this, let the component have $n$ vertices and $p$ original terminals.
Its degree sum is $3n-p-1$, so its rank is $(n-p+1)/2$.
Since $n\ge p$, this integer is positive.
Consequently, a bridge must separate $\{a,b\}$ from $c$.
Otherwise, one saturated terminal pair lies on one side and avoids a cycle on the other.

The endpoint of this bridge on the $\{a,b\}$ side cannot be $a$ or $b$.
Such a terminal path to $c$ would immediately cross the bridge and avoid a cycle on its starting side.
The other endpoint cannot be $c$, by the same argument.
Write the bridge as $zw$, with $a,b,z$ on its left side.
After deleting the bridge, the left graph $T_L$ is a three-terminal graph with terminals $a,b,z$.
The right graph $T_R$ has two degree-two terminals $w,c$ and otherwise degree three.
Write their ranks as $h_L,h_R$.
Both ranks are positive, and $h=h_L+h_R$.
Terminal paths factor across the bridge.
The bounds $P_{az}(T_L)\le2^{h_L}$ and $P_{wc}(T_R)\le2^{h_R}$ force equality in both.
The same argument gives $P_{bz}(T_L)=2^{h_L}$.
Also, $P_{ab}(T)=P_{ab}(T_L)$.
Induction and the monotonicity of $M$ prove the result when a bridge exists.

Now suppose that $T$ has no bridge.
A connected bridgeless graph of maximum degree three has no cut vertex.
Otherwise, each component after deleting a cut vertex requires at least two incident edges at that vertex.
Thus $T$ is two-connected.
When $h=1$, degree counting gives three vertices, so $T$ is a triangle.
Assume $h\ge2$.

The vertex $c$ is not adjacent to $a$ or $b$.
For example, an edge $ac$ gives at most
$1+2^{h-1}<2^h$ simple $a$--$c$ paths.
Here, delete that nonbridge edge to bound all paths other than the edge itself.
The two neighbors $x,y$ of $c$ are distinct.
If both incident edges joined $c$ to one vertex, that vertex's remaining edge would be a bridge.

Set $K=T-c$.
This connected four-terminal graph has terminals $a,b,x,y$ and rank $h-1$.
The equalities
\[
 P_{ax}(K)+P_{ay}(K)=2^h,
 \qquad
 P_{bx}(K)+P_{by}(K)=2^h
\]
force each summand to equal $2^{h-1}$.
Thus $K$ is saturated.
It has a through pair.
Otherwise, the vertex form of Menger's theorem gives a single vertex separating
$\{a,b\}$ from $\{x,y\}$ in $K$.
That vertex would disconnect $T$, contrary to two-connectivity.
The statement remains valid when the separating vertex is one of the four terminals.
At least one left terminal and one right terminal remain after its deletion.

Apply Proposition~\ref{prop:saturated-four-pole}.
Suppose that its ladder has $A$ square blocks and $s$ single-rung blocks.
Then $h=2A+s$, and its through-pair count is $2^A$.
An $a$--$b$ path in $K$ reaches one final block, turns there, and returns along the other strand.
A square block has two possible turns and two through choices.
A single-rung block has one turn and one through choice.
The square turns contribute $2(2^A-1)$ in total.
Each single-rung turn contributes at most $2^A$.
Therefore
\[
 P_{ab}(K)\le (s+2)2^A-2.
\]
Paths in $T$ through $c$ correspond exactly to the through pairs in $K$.
Hence
\[
 P_{ab}(T)=P_{ab}(K)+2^A
 \le(s+3)2^A-2.
\]
Subject to $2A+s=h$, this expression is largest when $A=\lfloor h/2\rfloor$.
Indeed, replacing two single-rung blocks by one square multiplies the positive factor by
$2(s+1)/(s+3)>1$ when $s\ge2$.
The resulting bound is $M(h)$.

For even $h$, take only square blocks and join the two right terminals to a new terminal $c$.
For odd $h\ge3$, add one single-rung block at the right end before adding $c$.
All single-rung turns then have their maximum weight, so equality holds.
For $h=1$, use a triangle.
All these examples are simple.
\end{proof}

\begin{lemma}\label{lem:saturated-two-pole-cycles}
Let $N$ be a connected loopless multigraph with two distinct degree-two terminals $u,v$.
Every other vertex has degree three.
If $\beta(N)=b$ and $P_{uv}(N)=2^b$, then
\[
 c(N)\le S(b):=\sum_{j=1}^b\bigl(M(j)-1\bigr)
 \le10\cdot2^{b/2}.
\]
The same bound holds if extra degree-two vertices subdivide edges.
\end{lemma}

\begin{proof}
Suppress unmarked vertices of degree two.
This preserves rank, terminal-path counts, and cycle counts.
No loop can arise in this normalization.
A pendant cycle would avoid every terminal path and contradict saturation.
The rank-one case is a cycle, with one cycle in total.

Suppose that $b\ge2$, and choose an edge $e=ux$ incident to $u$.
The vertices $u,v$ cannot be adjacent, since a direct edge would give at most $1+2^{b-1}$ terminal paths.
This edge is not a bridge.
If it were, both shores would contain cycles by degree counting.
If $v$ lies opposite $u$, every terminal path immediately crosses the bridge and avoids the cycle on the $u$ shore.
If $v$ lies on the $u$ shore, every terminal path avoids the cycle on the opposite shore.

Subdivide $e$ by a new vertex $a$ to obtain a three-terminal graph $T$ with terminals $u,a,v$.
Sliding the initial endpoint across the new edge gives a bijection between
$u$--$v$ paths and $a$--$v$ paths.
If a path starts with $ua$, remove that edge.
Otherwise, prepend $au$.
The degree-two condition at $u$ makes both operations reversible.
Thus both counts equal $2^b$.
Theorem~\ref{thm:saturated-three-pole} gives
\[
 1+c_e(N)=P_{ua}(T)\le M(b).
\]

Deleting $e$ leaves a connected graph of rank $b-1$.
Exactly half of the affine $u$--$v$ parity class avoids $e$.
Every member of that class is a path, so
$P_{uv}(N-e)=2^{b-1}$.
Delete the now-forced initial leaf segment and suppress unmarked degree-two vertices.
These operations preserve the rank and cycle count, and preserve the terminal-path count after moving the initial terminal.
If the rank is positive, the resulting terminals remain distinct and have degree two.
Otherwise, a forced terminal path would be unique, contrary to $2^{b-1}>1$.
A loop or an off-route cyclic component would contradict saturation.
Thus induction applies to the resulting two-terminal graph.
The identity $c(N)=c(N-e)+c_e(N)$ proves the recurrence.

Direct summation gives
\[
 S(2j)=10\cdot2^j-6j-10,
 \qquad
 S(2j+1)=14\cdot2^j-6j-13.
\]
Both expressions are at most $10\cdot2^{b/2}$.
\end{proof}

\begin{corollary}\label{cor:saturated-four-pole-cap}
For a saturated four-terminal graph $K$ of rank $h$, let $D(K)$ count its through pairs.
Then
\[
 \begin{split}
 c(K)&\le10\cdot2^{h/2},\\
 P_{L_1L_2}(K)+D(K)&\le M(h+1),\\
 P_{R_1R_2}(K)+D(K)&\le M(h+1),\\
 D(K)&\le2^{\lfloor(h+1)/2\rfloor}.
 \end{split}
\]
\end{corollary}

\begin{proof}
For the first bound, select one cross-terminal pair and suppress the two unused degree-two terminals.
Apply Lemma~\ref{lem:saturated-two-pole-cycles}.
For the second, add a new vertex $c$ joined to $R_1,R_2$.
The resulting three-terminal graph has rank $h+1$.
Each path count from $L_i$ to $c$ is $2^{h+1}$.
Its $L_1$--$L_2$ paths either avoid $c$ or correspond to through pairs in $K$.
Apply Theorem~\ref{thm:saturated-three-pole}.
The third bound is symmetric.
For the last bound, use Proposition~\ref{prop:saturated-four-pole} when $D(K)>0$.
When $D(K)=0$, it is immediate.
\end{proof}

\begin{lemma}\label{lem:saturated-four-pole-double-cap}
For a saturated four-terminal graph $K$ of rank $h$, let $C(K)$ count pairs
of vertex-disjoint cap paths, one from $L_1$ to $L_2$ and one from $R_1$ to $R_2$.
Then
\[
 C(K)\le3(h+1)2^{(h+1)/2}.
\]
\end{lemma}

\begin{proof}
First suppose that $D(K)>0$.
Use its classified ladder, and give a square block weight two and a single-rung block weight one.
Their total weight is $w=h+1$.
A disjoint cap pair turns its left path in a block $i$ and its right path in a later block.
The two paths cannot turn in the same block, since their local turns would intersect.
Fix the first turn block $i$.
The preceding through choices contribute at most $2^{w_{<i}/2}$.
The turn at block $i$ contributes at most $2^{w_i/2}$.
The right cap in the remaining suffix contributes at most $3\cdot2^{w_{>i}/2}$.
For this last bound, a suffix with $A$ squares and $s$ single rungs has at most $(s+2)2^A-2$ caps.
The inequality follows from $(s+2)2^{-s/2}\le3$.
It also covers a one-rung suffix, while an empty suffix contributes zero.
Thus each first-turn block contributes at most $3\cdot2^{w/2}$.
There are at most $w$ blocks.

Now suppose that $D(K)=0$.
Join $R_1,R_2$ to a new vertex $c$, as in Corollary~\ref{cor:saturated-four-pole-cap}.
The resulting saturated three-terminal graph has a bridge.
Otherwise, the bridgeless case in the proof of Theorem~\ref{thm:saturated-three-pole} would give a through pair in $K$.
Every such bridge separates $\{L_1,L_2\}$ from $c$, and neither endpoint is $c$.
It therefore lies in $K$ and separates its left terminal pair from its right terminal pair.

Each component of $K$ after deleting this bridge has positive rank.
Indeed, the same degree count gives rank $(n-p+1)/2\ge1$ for a side containing $p$ original terminals.
Neither bridge endpoint can be an original terminal.
Otherwise, a cross path from that terminal would avoid a cycle on its own side.
The two sides are thus three-terminal graphs, of ranks $h_L,h_R\ge1$, with $h_L+h_R=h$.
The cross-path counts factor across the bridge.
Their saturation forces both path counts to the bridge endpoint to be saturated on each side.
Hence Theorem~\ref{thm:saturated-three-pole} gives
\[
 C(K)=P_{L_1L_2}(K_L)P_{R_1R_2}(K_R)
 \le M(h_L)M(h_R)\le9\cdot2^{h/2}.
\]
Since $h\ge2$, this is at most $3(h+1)2^{(h+1)/2}$.
\end{proof}

\begin{theorem}\label{thm:saturated-ring-barrier}
Take $k\ge2$ vertex-disjoint saturated four-terminal graphs $K_1,\ldots,K_k$.
Join their right terminal pairs to the next left terminal pairs by two matching edges, cyclically.
The matching at each interface can be straight or crossed.
Let $G$ be the resulting graph and let $r=\beta(G)$.
Then
\[
 2^{r-1}\le c(G)\le2^{r-1}+20r2^{r/2}.
\]
Thus this entire saturated-block construction method cannot improve the half exponent in the cycle excess.
\end{theorem}

\begin{proof}
Write $h_i=\beta(K_i)$ and $w_i=h_i+1$.
Each four-terminal graph has $h_i\ge1$, so $w_i\ge2$.
Set $W=\sum_i w_i=r-1$.
Cut parity makes the number of used edges at every interface have the same parity.

For odd parity, choose one edge at each interface.
Each block then supplies $2^{h_i}$ paths between its selected cross terminals.
Every choice gives one simple cycle.
This sector has exactly
$2^k\prod_i2^{h_i}=2^W$ cycles.

An even-parity cycle with no used interface lies inside one block.
Corollary~\ref{cor:saturated-four-pole-cap} bounds their total by
$10\sum_i2^{h_i/2}\le10k2^{W/2}$.
If all interfaces are used, each block supplies a through pair or a pair of disjoint caps.
Two or more cap-pair blocks produce two or more components, so a single cycle uses at most one.
The all-through choices contribute at most $\prod_iD(K_i)\le2^{W/2}$ cycles.
By Lemma~\ref{lem:saturated-four-pole-double-cap}, the choices with one cap-pair block contribute at most
\[
 \sum_i C(K_i)\prod_{j\ne i}D(K_j)
 \le3W2^{W/2}.
\]

Otherwise, the used interfaces form a nonempty proper cyclic interval.
The two end blocks supply cap paths.
Every inner block supplies a through pair.
For a fixed interval whose occupied blocks form $I$, its contribution is at most
\[
 9\cdot2^{\frac12\sum_{i\in I}w_i}.
\]
Here we used $M(h_i+1)\le3\cdot2^{w_i/2}$ at both ends.
If the interval omits $t$ blocks, then their total weight is at least $2t$.
There are $k$ intervals for each $t\in\{0,\ldots,k-2\}$.
Their combined contribution is therefore at most
\[
 9k2^{W/2}\sum_{t=0}^{k-2}2^{-t}
 \le18k2^{W/2}.
\]
The total excess is at most $(28k+3W+1)2^{W/2}$.
Since $2k\le W$, this is at most $(17W+1)2^{W/2}$,
which is smaller than $20r2^{r/2}$.
\end{proof}
